% =====================================================================
%  BUỔI 5 — LỜI GIẢI BÀI TẬP VỀ NHÀ (từ giao-vien/02-loi-giai-bai-tap-ve-nha.md)
%  Sinh tự động bằng scripts/md2tex-luoigiai.py — không sửa tay.
%  Tài liệu giáo viên: dùng để chữa bài trên lớp và soạn đề kiểm tra.
% =====================================================================
\documentclass[11pt, a4paper]{article}
% =====================================================================
%  CAE SHANGHAI — GÓI LỆNH DÙNG CHUNG CHO TÀI LIỆU LUYỆN THI CSCA TOÁN
%  Biên dịch bằng XeLaTeX. Mọi file .tex trong buổi học \input file này.
% =====================================================================

% ---------- NGÔN NGỮ & PHÔNG CHỮ ----------
\usepackage{fontspec}
\usepackage{polyglossia}
\setmainlanguage{vietnamese}
\setotherlanguage{english}

% Chữ Latin (Anh + Việt): TeX Gyre Termes — bản clone Times New Roman,
% đi kèm TeX Live, hỗ trợ đầy đủ dấu tiếng Việt.
\setmainfont{TeX Gyre Termes}
\setsansfont{TeX Gyre Heros}

% Chữ Hán: ưu tiên Source Han Serif SC (思源宋体) — chuẩn học thuật Trung Quốc.
% Nếu máy chưa cài, tự động dùng Songti SC (宋体, có sẵn trên macOS).
% Cả hai đều thuộc họ 宋体 nên diện mạo tài liệu không đổi.
\IfFontExistsTF{Source Han Serif SC}{
  \newcommand{\cjkfamilyname}{Source Han Serif SC}
}{
  \IfFontExistsTF{Songti SC}{
    \newcommand{\cjkfamilyname}{Songti SC}
  }{
    \newcommand{\cjkfamilyname}{Noto Serif CJK SC}
  }
}

% xeCJK tự chuyển font khi gặp ký tự Hán, không cần bọc \cjk{} thủ công.
% Đây là điều kiện bắt buộc: font Latin (TeX Gyre Termes) không có glyph chữ Hán,
% nếu không khai báo xeCJK thì mọi chữ Hán viết trần sẽ bị nuốt khỏi bản in.
\usepackage{xeCJK}
\setCJKfamilyfont{caehan}{\cjkfamilyname}[Scale=0.95]
\setCJKmainfont{\cjkfamilyname}[Scale=0.95]
\setCJKsansfont{\cjkfamilyname}[Scale=0.95]
\setCJKmonofont{\cjkfamilyname}[Scale=0.95]

% \cjk{} ép dùng font Hán trong ngữ cảnh cần thiết (ví dụ trong tiêu đề in đậm
% hoặc trong bảng có định dạng riêng). Dùng \CJKfamily chứ không phải tên font,
% vì tên font chỉ là văn bản — sẽ in ra chữ thay vì đổi font.
\newcommand{\cjk}[1]{{\CJKfamily{caehan}#1}}

% ---------- TOÁN HỌC ----------
\usepackage{amsmath, amssymb, amsthm, mathtools}
\usepackage{unicode-math}
\setmathfont{TeX Gyre Termes Math}

% ---------- BỐ CỤC TRANG ----------
% Lề phải chừa vùng an toàn của letterhead:
%   banner trên chiếm 0–19mm  → top tối thiểu 2,2cm
%   banner dưới chiếm 281–297mm → bottom tối thiểu 2,2cm
% Lề trên 2,6cm (26mm) để tiêu đề đầu trang không đè banner (banner hết ở 19mm).
\usepackage[a4paper, top=2.6cm, bottom=2.4cm, left=2.5cm, right=2.0cm]{geometry}
\usepackage{fancyhdr}
\usepackage{titlesec}
\usepackage{enumitem}
\usepackage{graphicx}
\usepackage{booktabs, array, longtable, tabularx}
\usepackage{xcolor}
\usepackage[most]{tcolorbox}
\usepackage[colorlinks=true, linkcolor=cae_blue, citecolor=cae_blue, urlcolor=cae_blue]{hyperref}

% ---------- BẢNG MÀU THƯƠNG HIỆU ----------
\definecolor{cae_blue}{RGB}{0, 51, 102}
\definecolor{cae_gold}{RGB}{196, 149, 60}
\definecolor{box_bg}{RGB}{245, 247, 250}
\definecolor{box_line}{RGB}{0, 51, 102}

% ---------- TIÊU ĐỀ MỤC ----------
% Đệm đầu trang: banner letterhead chiếm 0–19mm, mà lề trên là 22mm.
% \section* đầu trang (sau \clearpage) sát lề nên phải chừa thêm, nếu không
% sẽ đè lên banner.
\titleformat{\section}
  {\normalfont\large\bfseries\sffamily\color{cae_blue}}{\thesection.}{0.6em}{}
\titleformat{\subsection}
  {\normalfont\normalsize\bfseries\sffamily\color{cae_blue!85!black}}{\thesubsection}{0.6em}{}
\titlespacing*{\section}{0pt}{14pt}{6pt}
\titlespacing*{\subsection}{0pt}{10pt}{4pt}

% Tiêu đề không đánh số ở đầu trang: \section* sau \clearpage nằm sát lề trên
% nên đè lên banner letterhead. Dùng \vspace* ngay TRƯỚC \section* không đủ,
% vì titlesec chèn khoảng cách riêng. Cách chắc chắn là đệm bằng \null.
\newcommand{\sectiondautrang}[1]{\section*{#1}}

% Dạng bài dùng \subsection*{} nên cần định dạng riêng, đậm và có gạch chân màu
\titleformat{\subsubsection}
  {\normalfont\normalsize\bfseries\sffamily\color{cae_blue!85!black}}{\thesubsubsection}{0.6em}{}

% ---------- LETTERHEAD (ảnh nền toàn trang) ----------
% Ảnh letterhead của CAE Shanghai: banner xanh trên/dưới kèm logo chìm giữa.
% Tỉ lệ ảnh 1414x2000 = 0.7070, khớp A4 (0.7071) nên chèn thẳng không méo.
% Đặt ở lớp nền dưới cùng để chữ không bị ảnh che.
% Dùng \AddToShipoutPictureBG (không có dấu *) để áp dụng cho MỌI trang.
\usepackage{eso-pic}
\AddToShipoutPictureBG{%
  \AtPageLowerLeft{%
    \includegraphics[width=\paperwidth,height=\paperheight]{images/letterhead.png}%
  }%
}

% ---------- ĐẦU / CHÂN TRANG ----------
% Letterhead đã có banner trên và dưới, nên bỏ dòng chạy để tránh chồng chữ.
% Số trang đặt trong vùng nội dung, sát trên banner dưới, lề phải.
\pagestyle{fancy}
\setlength{\headheight}{0pt}
\fancyhf{}
\renewcommand{\headrulewidth}{0pt}
\renewcommand{\footrulewidth}{0pt}
\fancyfoot[R]{\raisebox{0.9cm}{\small\color{cae_blue}\thepage}}

% ---------- MÔI TRƯỜNG NỘI DUNG ----------
% Trong tcolorbox, LaTeX đặt parindent = parskip = 0, nên các đoạn văn trong
% khung dính liền nhau. Đặt lại khoảng cách đoạn cho mọi khung.
\newcommand{\khoangcachdoan}{\setlength{\parskip}{0.5em}\setlength{\parindent}{0pt}}
\tcbset{before upper={\khoangcachdoan}}

\newtcolorbox{dinhnghia}[1][]{
  colback=box_bg, colframe=box_line, boxrule=0.6pt, arc=1pt,
  fonttitle=\bfseries, title={Định nghĩa #1}, breakable}
\newtcolorbox{dinhly}[1][]{
  colback=box_bg, colframe=cae_gold, boxrule=0.6pt, arc=1pt,
  fonttitle=\bfseries, title={Tính chất #1}, breakable}
\newtcolorbox{vidu}[1][]{
  colback=white, colframe=black!40, boxrule=0.4pt, arc=1pt,
  fonttitle=\bfseries, title={Ví dụ #1}, breakable}
\newtcolorbox{baitap}[1][]{
  colback=white, colframe=cae_blue!50, boxrule=0.4pt, arc=1pt,
  fonttitle=\bfseries, title={Câu #1}, breakable}

% Ô đáp án: viền xanh đậm hơn, dùng trong file đáp án
\newtcolorbox{dapan}[1][]{
  colback=box_bg, colframe=cae_blue, boxrule=0.6pt, arc=1pt,
  fonttitle=\bfseries, title={Câu #1}, breakable}

% ---------- TÊN CHỨNG MINH ----------
\renewcommand{\proofname}{\normalfont\bfseries Chứng minh}

% ---------- MÔN HỌC ----------
% Đổi được ở file .tex trước khi % =====================================================================
%  CAE SHANGHAI — GÓI LỆNH DÙNG CHUNG CHO TÀI LIỆU LUYỆN THI CSCA TOÁN
%  Biên dịch bằng XeLaTeX. Mọi file .tex trong buổi học \input file này.
% =====================================================================

% ---------- NGÔN NGỮ & PHÔNG CHỮ ----------
\usepackage{fontspec}
\usepackage{polyglossia}
\setmainlanguage{vietnamese}
\setotherlanguage{english}

% Chữ Latin (Anh + Việt): TeX Gyre Termes — bản clone Times New Roman,
% đi kèm TeX Live, hỗ trợ đầy đủ dấu tiếng Việt.
\setmainfont{TeX Gyre Termes}
\setsansfont{TeX Gyre Heros}

% Chữ Hán: ưu tiên Source Han Serif SC (思源宋体) — chuẩn học thuật Trung Quốc.
% Nếu máy chưa cài, tự động dùng Songti SC (宋体, có sẵn trên macOS).
% Cả hai đều thuộc họ 宋体 nên diện mạo tài liệu không đổi.
\IfFontExistsTF{Source Han Serif SC}{
  \newcommand{\cjkfamilyname}{Source Han Serif SC}
}{
  \IfFontExistsTF{Songti SC}{
    \newcommand{\cjkfamilyname}{Songti SC}
  }{
    \newcommand{\cjkfamilyname}{Noto Serif CJK SC}
  }
}

% xeCJK tự chuyển font khi gặp ký tự Hán, không cần bọc \cjk{} thủ công.
% Đây là điều kiện bắt buộc: font Latin (TeX Gyre Termes) không có glyph chữ Hán,
% nếu không khai báo xeCJK thì mọi chữ Hán viết trần sẽ bị nuốt khỏi bản in.
\usepackage{xeCJK}
\setCJKfamilyfont{caehan}{\cjkfamilyname}[Scale=0.95]
\setCJKmainfont{\cjkfamilyname}[Scale=0.95]
\setCJKsansfont{\cjkfamilyname}[Scale=0.95]
\setCJKmonofont{\cjkfamilyname}[Scale=0.95]

% \cjk{} ép dùng font Hán trong ngữ cảnh cần thiết (ví dụ trong tiêu đề in đậm
% hoặc trong bảng có định dạng riêng). Dùng \CJKfamily chứ không phải tên font,
% vì tên font chỉ là văn bản — sẽ in ra chữ thay vì đổi font.
\newcommand{\cjk}[1]{{\CJKfamily{caehan}#1}}

% ---------- TOÁN HỌC ----------
\usepackage{amsmath, amssymb, amsthm, mathtools}
\usepackage{unicode-math}
\setmathfont{TeX Gyre Termes Math}

% ---------- BỐ CỤC TRANG ----------
% Lề phải chừa vùng an toàn của letterhead:
%   banner trên chiếm 0–19mm  → top tối thiểu 2,2cm
%   banner dưới chiếm 281–297mm → bottom tối thiểu 2,2cm
% Lề trên 2,6cm (26mm) để tiêu đề đầu trang không đè banner (banner hết ở 19mm).
\usepackage[a4paper, top=2.6cm, bottom=2.4cm, left=2.5cm, right=2.0cm]{geometry}
\usepackage{fancyhdr}
\usepackage{titlesec}
\usepackage{enumitem}
\usepackage{graphicx}
\usepackage{booktabs, array, longtable, tabularx}
\usepackage{xcolor}
\usepackage[most]{tcolorbox}
\usepackage[colorlinks=true, linkcolor=cae_blue, citecolor=cae_blue, urlcolor=cae_blue]{hyperref}

% ---------- BẢNG MÀU THƯƠNG HIỆU ----------
\definecolor{cae_blue}{RGB}{0, 51, 102}
\definecolor{cae_gold}{RGB}{196, 149, 60}
\definecolor{box_bg}{RGB}{245, 247, 250}
\definecolor{box_line}{RGB}{0, 51, 102}

% ---------- TIÊU ĐỀ MỤC ----------
% Đệm đầu trang: banner letterhead chiếm 0–19mm, mà lề trên là 22mm.
% \section* đầu trang (sau \clearpage) sát lề nên phải chừa thêm, nếu không
% sẽ đè lên banner.
\titleformat{\section}
  {\normalfont\large\bfseries\sffamily\color{cae_blue}}{\thesection.}{0.6em}{}
\titleformat{\subsection}
  {\normalfont\normalsize\bfseries\sffamily\color{cae_blue!85!black}}{\thesubsection}{0.6em}{}
\titlespacing*{\section}{0pt}{14pt}{6pt}
\titlespacing*{\subsection}{0pt}{10pt}{4pt}

% Tiêu đề không đánh số ở đầu trang: \section* sau \clearpage nằm sát lề trên
% nên đè lên banner letterhead. Dùng \vspace* ngay TRƯỚC \section* không đủ,
% vì titlesec chèn khoảng cách riêng. Cách chắc chắn là đệm bằng \null.
\newcommand{\sectiondautrang}[1]{\section*{#1}}

% Dạng bài dùng \subsection*{} nên cần định dạng riêng, đậm và có gạch chân màu
\titleformat{\subsubsection}
  {\normalfont\normalsize\bfseries\sffamily\color{cae_blue!85!black}}{\thesubsubsection}{0.6em}{}

% ---------- LETTERHEAD (ảnh nền toàn trang) ----------
% Ảnh letterhead của CAE Shanghai: banner xanh trên/dưới kèm logo chìm giữa.
% Tỉ lệ ảnh 1414x2000 = 0.7070, khớp A4 (0.7071) nên chèn thẳng không méo.
% Đặt ở lớp nền dưới cùng để chữ không bị ảnh che.
% Dùng \AddToShipoutPictureBG (không có dấu *) để áp dụng cho MỌI trang.
\usepackage{eso-pic}
\AddToShipoutPictureBG{%
  \AtPageLowerLeft{%
    \includegraphics[width=\paperwidth,height=\paperheight]{images/letterhead.png}%
  }%
}

% ---------- ĐẦU / CHÂN TRANG ----------
% Letterhead đã có banner trên và dưới, nên bỏ dòng chạy để tránh chồng chữ.
% Số trang đặt trong vùng nội dung, sát trên banner dưới, lề phải.
\pagestyle{fancy}
\setlength{\headheight}{0pt}
\fancyhf{}
\renewcommand{\headrulewidth}{0pt}
\renewcommand{\footrulewidth}{0pt}
\fancyfoot[R]{\raisebox{0.9cm}{\small\color{cae_blue}\thepage}}

% ---------- MÔI TRƯỜNG NỘI DUNG ----------
% Trong tcolorbox, LaTeX đặt parindent = parskip = 0, nên các đoạn văn trong
% khung dính liền nhau. Đặt lại khoảng cách đoạn cho mọi khung.
\newcommand{\khoangcachdoan}{\setlength{\parskip}{0.5em}\setlength{\parindent}{0pt}}
\tcbset{before upper={\khoangcachdoan}}

\newtcolorbox{dinhnghia}[1][]{
  colback=box_bg, colframe=box_line, boxrule=0.6pt, arc=1pt,
  fonttitle=\bfseries, title={Định nghĩa #1}, breakable}
\newtcolorbox{dinhly}[1][]{
  colback=box_bg, colframe=cae_gold, boxrule=0.6pt, arc=1pt,
  fonttitle=\bfseries, title={Tính chất #1}, breakable}
\newtcolorbox{vidu}[1][]{
  colback=white, colframe=black!40, boxrule=0.4pt, arc=1pt,
  fonttitle=\bfseries, title={Ví dụ #1}, breakable}
\newtcolorbox{baitap}[1][]{
  colback=white, colframe=cae_blue!50, boxrule=0.4pt, arc=1pt,
  fonttitle=\bfseries, title={Câu #1}, breakable}

% Ô đáp án: viền xanh đậm hơn, dùng trong file đáp án
\newtcolorbox{dapan}[1][]{
  colback=box_bg, colframe=cae_blue, boxrule=0.6pt, arc=1pt,
  fonttitle=\bfseries, title={Câu #1}, breakable}

% ---------- TÊN CHỨNG MINH ----------
\renewcommand{\proofname}{\normalfont\bfseries Chứng minh}

% ---------- MÔN HỌC ----------
% Đổi được ở file .tex trước khi % =====================================================================
%  CAE SHANGHAI — GÓI LỆNH DÙNG CHUNG CHO TÀI LIỆU LUYỆN THI CSCA TOÁN
%  Biên dịch bằng XeLaTeX. Mọi file .tex trong buổi học \input file này.
% =====================================================================

% ---------- NGÔN NGỮ & PHÔNG CHỮ ----------
\usepackage{fontspec}
\usepackage{polyglossia}
\setmainlanguage{vietnamese}
\setotherlanguage{english}

% Chữ Latin (Anh + Việt): TeX Gyre Termes — bản clone Times New Roman,
% đi kèm TeX Live, hỗ trợ đầy đủ dấu tiếng Việt.
\setmainfont{TeX Gyre Termes}
\setsansfont{TeX Gyre Heros}

% Chữ Hán: ưu tiên Source Han Serif SC (思源宋体) — chuẩn học thuật Trung Quốc.
% Nếu máy chưa cài, tự động dùng Songti SC (宋体, có sẵn trên macOS).
% Cả hai đều thuộc họ 宋体 nên diện mạo tài liệu không đổi.
\IfFontExistsTF{Source Han Serif SC}{
  \newcommand{\cjkfamilyname}{Source Han Serif SC}
}{
  \IfFontExistsTF{Songti SC}{
    \newcommand{\cjkfamilyname}{Songti SC}
  }{
    \newcommand{\cjkfamilyname}{Noto Serif CJK SC}
  }
}

% xeCJK tự chuyển font khi gặp ký tự Hán, không cần bọc \cjk{} thủ công.
% Đây là điều kiện bắt buộc: font Latin (TeX Gyre Termes) không có glyph chữ Hán,
% nếu không khai báo xeCJK thì mọi chữ Hán viết trần sẽ bị nuốt khỏi bản in.
\usepackage{xeCJK}
\setCJKfamilyfont{caehan}{\cjkfamilyname}[Scale=0.95]
\setCJKmainfont{\cjkfamilyname}[Scale=0.95]
\setCJKsansfont{\cjkfamilyname}[Scale=0.95]
\setCJKmonofont{\cjkfamilyname}[Scale=0.95]

% \cjk{} ép dùng font Hán trong ngữ cảnh cần thiết (ví dụ trong tiêu đề in đậm
% hoặc trong bảng có định dạng riêng). Dùng \CJKfamily chứ không phải tên font,
% vì tên font chỉ là văn bản — sẽ in ra chữ thay vì đổi font.
\newcommand{\cjk}[1]{{\CJKfamily{caehan}#1}}

% ---------- TOÁN HỌC ----------
\usepackage{amsmath, amssymb, amsthm, mathtools}
\usepackage{unicode-math}
\setmathfont{TeX Gyre Termes Math}

% ---------- BỐ CỤC TRANG ----------
% Lề phải chừa vùng an toàn của letterhead:
%   banner trên chiếm 0–19mm  → top tối thiểu 2,2cm
%   banner dưới chiếm 281–297mm → bottom tối thiểu 2,2cm
% Lề trên 2,6cm (26mm) để tiêu đề đầu trang không đè banner (banner hết ở 19mm).
\usepackage[a4paper, top=2.6cm, bottom=2.4cm, left=2.5cm, right=2.0cm]{geometry}
\usepackage{fancyhdr}
\usepackage{titlesec}
\usepackage{enumitem}
\usepackage{graphicx}
\usepackage{booktabs, array, longtable, tabularx}
\usepackage{xcolor}
\usepackage[most]{tcolorbox}
\usepackage[colorlinks=true, linkcolor=cae_blue, citecolor=cae_blue, urlcolor=cae_blue]{hyperref}

% ---------- BẢNG MÀU THƯƠNG HIỆU ----------
\definecolor{cae_blue}{RGB}{0, 51, 102}
\definecolor{cae_gold}{RGB}{196, 149, 60}
\definecolor{box_bg}{RGB}{245, 247, 250}
\definecolor{box_line}{RGB}{0, 51, 102}

% ---------- TIÊU ĐỀ MỤC ----------
% Đệm đầu trang: banner letterhead chiếm 0–19mm, mà lề trên là 22mm.
% \section* đầu trang (sau \clearpage) sát lề nên phải chừa thêm, nếu không
% sẽ đè lên banner.
\titleformat{\section}
  {\normalfont\large\bfseries\sffamily\color{cae_blue}}{\thesection.}{0.6em}{}
\titleformat{\subsection}
  {\normalfont\normalsize\bfseries\sffamily\color{cae_blue!85!black}}{\thesubsection}{0.6em}{}
\titlespacing*{\section}{0pt}{14pt}{6pt}
\titlespacing*{\subsection}{0pt}{10pt}{4pt}

% Tiêu đề không đánh số ở đầu trang: \section* sau \clearpage nằm sát lề trên
% nên đè lên banner letterhead. Dùng \vspace* ngay TRƯỚC \section* không đủ,
% vì titlesec chèn khoảng cách riêng. Cách chắc chắn là đệm bằng \null.
\newcommand{\sectiondautrang}[1]{\section*{#1}}

% Dạng bài dùng \subsection*{} nên cần định dạng riêng, đậm và có gạch chân màu
\titleformat{\subsubsection}
  {\normalfont\normalsize\bfseries\sffamily\color{cae_blue!85!black}}{\thesubsubsection}{0.6em}{}

% ---------- LETTERHEAD (ảnh nền toàn trang) ----------
% Ảnh letterhead của CAE Shanghai: banner xanh trên/dưới kèm logo chìm giữa.
% Tỉ lệ ảnh 1414x2000 = 0.7070, khớp A4 (0.7071) nên chèn thẳng không méo.
% Đặt ở lớp nền dưới cùng để chữ không bị ảnh che.
% Dùng \AddToShipoutPictureBG (không có dấu *) để áp dụng cho MỌI trang.
\usepackage{eso-pic}
\AddToShipoutPictureBG{%
  \AtPageLowerLeft{%
    \includegraphics[width=\paperwidth,height=\paperheight]{images/letterhead.png}%
  }%
}

% ---------- ĐẦU / CHÂN TRANG ----------
% Letterhead đã có banner trên và dưới, nên bỏ dòng chạy để tránh chồng chữ.
% Số trang đặt trong vùng nội dung, sát trên banner dưới, lề phải.
\pagestyle{fancy}
\setlength{\headheight}{0pt}
\fancyhf{}
\renewcommand{\headrulewidth}{0pt}
\renewcommand{\footrulewidth}{0pt}
\fancyfoot[R]{\raisebox{0.9cm}{\small\color{cae_blue}\thepage}}

% ---------- MÔI TRƯỜNG NỘI DUNG ----------
% Trong tcolorbox, LaTeX đặt parindent = parskip = 0, nên các đoạn văn trong
% khung dính liền nhau. Đặt lại khoảng cách đoạn cho mọi khung.
\newcommand{\khoangcachdoan}{\setlength{\parskip}{0.5em}\setlength{\parindent}{0pt}}
\tcbset{before upper={\khoangcachdoan}}

\newtcolorbox{dinhnghia}[1][]{
  colback=box_bg, colframe=box_line, boxrule=0.6pt, arc=1pt,
  fonttitle=\bfseries, title={Định nghĩa #1}, breakable}
\newtcolorbox{dinhly}[1][]{
  colback=box_bg, colframe=cae_gold, boxrule=0.6pt, arc=1pt,
  fonttitle=\bfseries, title={Tính chất #1}, breakable}
\newtcolorbox{vidu}[1][]{
  colback=white, colframe=black!40, boxrule=0.4pt, arc=1pt,
  fonttitle=\bfseries, title={Ví dụ #1}, breakable}
\newtcolorbox{baitap}[1][]{
  colback=white, colframe=cae_blue!50, boxrule=0.4pt, arc=1pt,
  fonttitle=\bfseries, title={Câu #1}, breakable}

% Ô đáp án: viền xanh đậm hơn, dùng trong file đáp án
\newtcolorbox{dapan}[1][]{
  colback=box_bg, colframe=cae_blue, boxrule=0.6pt, arc=1pt,
  fonttitle=\bfseries, title={Câu #1}, breakable}

% ---------- TÊN CHỨNG MINH ----------
\renewcommand{\proofname}{\normalfont\bfseries Chứng minh}

% ---------- MÔN HỌC ----------
% Đổi được ở file .tex trước khi \input{cae-style} bằng \renewcommand{\monhoc}{...}
\providecommand{\monhoc}{Môn Toán}

% ---------- TRANG BÌA GỌN ----------
% #1 = loại tài liệu (TÀI LIỆU CHUẨN BỊ, TÀI LIỆU BUỔI HỌC, ...)
% #2 = số buổi
% #3 = tên buổi tiếng Việt
% #4 = tên buổi tiếng Hán
% #5 = module và tỉ trọng
\newcommand{\trangbia}[5]{%
  \begin{center}
    \vspace*{0.6cm}
    {\large\bfseries\color{cae_blue} CAE SHANGHAI}\\[0.15cm]
    {\footnotesize Tài liệu luyện thi CSCA · \monhoc}\\[0.5cm]
    {\color{cae_gold}\rule{0.55\linewidth}{0.8pt}}\\[0.35cm]
    {\small\bfseries\color{cae_blue!80!black} #1}\\[0.35cm]
    {\Large\bfseries\color{cae_blue} BUỔI #2 — #3}\\[0.2cm]
    {\large\color{cae_blue!85!black} \cjk{#4}}\\[0.25cm]
    \ifx\relax#5\relax\else{\small #5}\\[0.35cm]\fi
    {\color{cae_gold}\rule{0.55\linewidth}{0.8pt}}\\[0.3cm]
    {\footnotesize Biên soạn: Hoàng Dương, CAE Shanghai \quad·\quad \today}
  \end{center}
  \vspace{0.3cm}
}

% ---------- TRANG BÌA CHO ĐỀ KIỂM TRA ----------
% Khác \trangbia ở chỗ không có chữ "BUỔI N" — dùng cho đề thi, không phải bài giảng.
% #1 = loại tài liệu, #2 = tên kỳ thi (Việt), #3 = tên kỳ thi (Hán), #4 = thông số
\newcommand{\biade}[4]{%
  \begin{center}
    \vspace*{0.6cm}
    {\large\bfseries\color{cae_blue} CAE SHANGHAI}\\[0.15cm]
    {\footnotesize Tài liệu luyện thi CSCA · \monhoc}\\[0.5cm]
    {\color{cae_gold}\rule{0.55\linewidth}{0.8pt}}\\[0.35cm]
    {\small\bfseries\color{cae_blue!80!black} #1}\\[0.35cm]
    {\Large\bfseries\color{cae_blue} #2}\\[0.2cm]
    {\large\color{cae_blue!85!black} \cjk{#3}}\\[0.25cm]
    {\small #4}\\[0.35cm]
    {\color{cae_gold}\rule{0.55\linewidth}{0.8pt}}\\[0.3cm]
    {\footnotesize Biên soạn: Hoàng Dương, CAE Shanghai \quad·\quad \today}
  \end{center}
  \vspace{0.3cm}
}

% ---------- TIỆN ÍCH ----------
% Nhãn từ khóa in đậm ở đầu đoạn, dùng cho Từ vựng / Từ khóa nhận dạng / Chú ý
% Khoảng cách đoạn do khung tự đặt (xem \khoangcachdoan ở trên),
% ở đây chỉ cần tách đoạn và bỏ thụt đầu dòng.
\newcommand{\nhan}[1]{\par\noindent\textbf{#1.}\ \ignorespaces}

% Bảng gọn, không đánh số
\newcommand{\bang}[1]{%
  \begin{center}#1\end{center}}

% Đánh dấu ô trống cho học sinh điền
\newcommand{\chotrong}{\makebox[3.2cm]{\dotfill}}
 bằng \renewcommand{\monhoc}{...}
\providecommand{\monhoc}{Môn Toán}

% ---------- TRANG BÌA GỌN ----------
% #1 = loại tài liệu (TÀI LIỆU CHUẨN BỊ, TÀI LIỆU BUỔI HỌC, ...)
% #2 = số buổi
% #3 = tên buổi tiếng Việt
% #4 = tên buổi tiếng Hán
% #5 = module và tỉ trọng
\newcommand{\trangbia}[5]{%
  \begin{center}
    \vspace*{0.6cm}
    {\large\bfseries\color{cae_blue} CAE SHANGHAI}\\[0.15cm]
    {\footnotesize Tài liệu luyện thi CSCA · \monhoc}\\[0.5cm]
    {\color{cae_gold}\rule{0.55\linewidth}{0.8pt}}\\[0.35cm]
    {\small\bfseries\color{cae_blue!80!black} #1}\\[0.35cm]
    {\Large\bfseries\color{cae_blue} BUỔI #2 — #3}\\[0.2cm]
    {\large\color{cae_blue!85!black} \cjk{#4}}\\[0.25cm]
    \ifx\relax#5\relax\else{\small #5}\\[0.35cm]\fi
    {\color{cae_gold}\rule{0.55\linewidth}{0.8pt}}\\[0.3cm]
    {\footnotesize Biên soạn: Hoàng Dương, CAE Shanghai \quad·\quad \today}
  \end{center}
  \vspace{0.3cm}
}

% ---------- TRANG BÌA CHO ĐỀ KIỂM TRA ----------
% Khác \trangbia ở chỗ không có chữ "BUỔI N" — dùng cho đề thi, không phải bài giảng.
% #1 = loại tài liệu, #2 = tên kỳ thi (Việt), #3 = tên kỳ thi (Hán), #4 = thông số
\newcommand{\biade}[4]{%
  \begin{center}
    \vspace*{0.6cm}
    {\large\bfseries\color{cae_blue} CAE SHANGHAI}\\[0.15cm]
    {\footnotesize Tài liệu luyện thi CSCA · \monhoc}\\[0.5cm]
    {\color{cae_gold}\rule{0.55\linewidth}{0.8pt}}\\[0.35cm]
    {\small\bfseries\color{cae_blue!80!black} #1}\\[0.35cm]
    {\Large\bfseries\color{cae_blue} #2}\\[0.2cm]
    {\large\color{cae_blue!85!black} \cjk{#3}}\\[0.25cm]
    {\small #4}\\[0.35cm]
    {\color{cae_gold}\rule{0.55\linewidth}{0.8pt}}\\[0.3cm]
    {\footnotesize Biên soạn: Hoàng Dương, CAE Shanghai \quad·\quad \today}
  \end{center}
  \vspace{0.3cm}
}

% ---------- TIỆN ÍCH ----------
% Nhãn từ khóa in đậm ở đầu đoạn, dùng cho Từ vựng / Từ khóa nhận dạng / Chú ý
% Khoảng cách đoạn do khung tự đặt (xem \khoangcachdoan ở trên),
% ở đây chỉ cần tách đoạn và bỏ thụt đầu dòng.
\newcommand{\nhan}[1]{\par\noindent\textbf{#1.}\ \ignorespaces}

% Bảng gọn, không đánh số
\newcommand{\bang}[1]{%
  \begin{center}#1\end{center}}

% Đánh dấu ô trống cho học sinh điền
\newcommand{\chotrong}{\makebox[3.2cm]{\dotfill}}
 bằng \renewcommand{\monhoc}{...}
\providecommand{\monhoc}{Môn Toán}

% ---------- TRANG BÌA GỌN ----------
% #1 = loại tài liệu (TÀI LIỆU CHUẨN BỊ, TÀI LIỆU BUỔI HỌC, ...)
% #2 = số buổi
% #3 = tên buổi tiếng Việt
% #4 = tên buổi tiếng Hán
% #5 = module và tỉ trọng
\newcommand{\trangbia}[5]{%
  \begin{center}
    \vspace*{0.6cm}
    {\large\bfseries\color{cae_blue} CAE SHANGHAI}\\[0.15cm]
    {\footnotesize Tài liệu luyện thi CSCA · \monhoc}\\[0.5cm]
    {\color{cae_gold}\rule{0.55\linewidth}{0.8pt}}\\[0.35cm]
    {\small\bfseries\color{cae_blue!80!black} #1}\\[0.35cm]
    {\Large\bfseries\color{cae_blue} BUỔI #2 — #3}\\[0.2cm]
    {\large\color{cae_blue!85!black} \cjk{#4}}\\[0.25cm]
    \ifx\relax#5\relax\else{\small #5}\\[0.35cm]\fi
    {\color{cae_gold}\rule{0.55\linewidth}{0.8pt}}\\[0.3cm]
    {\footnotesize Biên soạn: Hoàng Dương, CAE Shanghai \quad·\quad \today}
  \end{center}
  \vspace{0.3cm}
}

% ---------- TRANG BÌA CHO ĐỀ KIỂM TRA ----------
% Khác \trangbia ở chỗ không có chữ "BUỔI N" — dùng cho đề thi, không phải bài giảng.
% #1 = loại tài liệu, #2 = tên kỳ thi (Việt), #3 = tên kỳ thi (Hán), #4 = thông số
\newcommand{\biade}[4]{%
  \begin{center}
    \vspace*{0.6cm}
    {\large\bfseries\color{cae_blue} CAE SHANGHAI}\\[0.15cm]
    {\footnotesize Tài liệu luyện thi CSCA · \monhoc}\\[0.5cm]
    {\color{cae_gold}\rule{0.55\linewidth}{0.8pt}}\\[0.35cm]
    {\small\bfseries\color{cae_blue!80!black} #1}\\[0.35cm]
    {\Large\bfseries\color{cae_blue} #2}\\[0.2cm]
    {\large\color{cae_blue!85!black} \cjk{#3}}\\[0.25cm]
    {\small #4}\\[0.35cm]
    {\color{cae_gold}\rule{0.55\linewidth}{0.8pt}}\\[0.3cm]
    {\footnotesize Biên soạn: Hoàng Dương, CAE Shanghai \quad·\quad \today}
  \end{center}
  \vspace{0.3cm}
}

% ---------- TIỆN ÍCH ----------
% Nhãn từ khóa in đậm ở đầu đoạn, dùng cho Từ vựng / Từ khóa nhận dạng / Chú ý
% Khoảng cách đoạn do khung tự đặt (xem \khoangcachdoan ở trên),
% ở đây chỉ cần tách đoạn và bỏ thụt đầu dòng.
\newcommand{\nhan}[1]{\par\noindent\textbf{#1.}\ \ignorespaces}

% Bảng gọn, không đánh số
\newcommand{\bang}[1]{%
  \begin{center}#1\end{center}}

% Đánh dấu ô trống cho học sinh điền
\newcommand{\chotrong}{\makebox[3.2cm]{\dotfill}}


\begin{document}

\trangbia{LỜI GIẢI BÀI TẬP VỀ NHÀ}{5}{Lượng giác: hàm số và phương trình}{三角函数的图像与方程}{Tài liệu giáo viên \quad·\quad 40 câu}

\subsection*{Nhóm A — Hàm số và phương trình lượng giác}

\begin{dapan}[1 — Đáp án A]
Hàm $\tan$ không xác định khi biểu thức bên trong bằng $\dfrac{\pi}{2} + k\pi$. Đặt $x + \dfrac{\pi}{4} = \dfrac{\pi}{2} + k\pi$, suy ra $x = \dfrac{\pi}{4} + k\pi$. Vậy tập xác định là $\left\{x \mid x \neq \dfrac{\pi}{4} + k\pi,\ k \in \mathbb{Z}\right\}$.

\textbf{Bẫy.} Dấu trong ngoặc là \textbf{cộng} $\dfrac{\pi}{4}$ nên nghiệm bị loại là $\dfrac{\pi}{4} + k\pi$, không phải $-\dfrac{\pi}{4} + k\pi$. Đảo dấu là sai.
\end{dapan}

\begin{dapan}[2 — Đáp án A]
Hàm $\sin x$ đồng biến trên $\left[-\dfrac{\pi}{2}, \dfrac{\pi}{2}\right]$. Khoảng $\left[0, \dfrac{\pi}{2}\right]$ nằm trọn trong khoảng đồng biến đó nên hàm số tăng trên toàn khoảng.
\end{dapan}

\begin{dapan}[3 — Đáp án A]
Hệ số của $x$ là $\omega = -1$, lấy trị tuyệt đối được $|\omega| = 1$. Chu kỳ là $T = \dfrac{2\pi}{1} = 2\pi$.

\textbf{Chú ý.} Dấu trừ của $\omega$ không ảnh hưởng đến chu kỳ vì công thức dùng $|\omega|$.
\end{dapan}

\begin{dapan}[4 — Đáp án B]
Phương trình $\sin x = \dfrac{\sqrt{3}}{2}$ trên $\left[0, \dfrac{\pi}{2}\right]$. Vì $\sin\dfrac{\pi}{3} = \dfrac{\sqrt{3}}{2}$ nên $x = \dfrac{\pi}{3}$.

\textbf{Chú ý.} Trên cả vòng tròn có hai nghiệm $\dfrac{\pi}{3}$ và $\dfrac{2\pi}{3}$, nhưng đề giới hạn trong $\left[0, \dfrac{\pi}{2}\right]$ nên chỉ lấy $\dfrac{\pi}{3}$.
\end{dapan}

\begin{dapan}[5 — Đáp án A]
Hàm $\tan x$ đồng biến trên từng khoảng $\left(-\dfrac{\pi}{2} + k\pi, \dfrac{\pi}{2} + k\pi\right)$. Khoảng $\left(\dfrac{\pi}{2}, \dfrac{3\pi}{2}\right)$ ứng với $k = 1$, nằm trọn trong một khoảng xác định nên hàm số tăng.
\end{dapan}

\begin{dapan}[6 — Đáp án C]
Biểu thức $y = \dfrac{1}{3}\sin 3x$ có $\omega = 3$. Tần số là $f = \dfrac{\omega}{2\pi} = \dfrac{3}{2\pi}$.

\textbf{Chú ý.} Phân biệt tần số $f = \dfrac{\omega}{2\pi}$ và chu kỳ $T = \dfrac{2\pi}{\omega}$. Đề hỏi 频率 (tần số) nên lấy $\omega$ chia $2\pi$.
\end{dapan}

\begin{dapan}[7 — Đáp án B]
Phương trình $\cos x = -1$ có nghiệm $x = \pi + 2k\pi$ với $k \in \mathbb{Z}$.
\end{dapan}

\begin{dapan}[8 — Đáp án C]
Từ $2\cos x + 1 = 0$ suy ra $\cos x = -\dfrac{1}{2}$. Trên $[0, 2\pi]$ có hai nghiệm $x = \dfrac{2\pi}{3}$ và $x = \dfrac{4\pi}{3}$.
\end{dapan}

\begin{dapan}[9 — Đáp án C]
Phương trình $\cos x = \dfrac{1}{2}$ trên $[0, \pi]$ có nghiệm duy nhất $x = \dfrac{\pi}{3}$.

\textbf{Chú ý.} Họ nghiệm tổng quát là $x = \pm\dfrac{\pi}{3} + 2k\pi$, nhưng nghiệm $x = -\dfrac{\pi}{3}$ nằm ngoài đoạn $[0, \pi]$ nên bị loại.
\end{dapan}

\begin{dapan}[10 — Đáp án C]
Đặt $t = \sin x$ với $t \in [-1, 1]$. Hàm số trở thành $y = 2t^2 - 4t + 1$, đỉnh tại $t_0 = \dfrac{4}{4} = 1$.

Vì $t_0 = 1$ nằm trong đoạn $[-1, 1]$ nên giá trị nhỏ nhất đạt tại $t = 1$: $y = 2 - 4 + 1 = -1$. Giá trị lớn nhất đạt tại đầu mút $t = -1$: $y = 2 + 4 + 1 = 7$.

Vậy tập giá trị là $[-1, 7]$.

\textbf{Chú ý.} Phải so sánh đỉnh $t_0$ với hai đầu mút $-1$ và $1$. Nếu thay $t_0$ một cách máy móc mà không kiểm tra thì dễ sai.
\end{dapan}

\begin{dapan}[11 — Đáp án C]
Hàm $\sin u$ đồng biến khi $-\dfrac{\pi}{2} + 2k\pi \le u \le \dfrac{\pi}{2} + 2k\pi$. Đặt $u = 2x - \dfrac{\pi}{3}$:

$$-\frac{\pi}{2} + 2k\pi \le 2x - \frac{\pi}{3} \le \frac{\pi}{2} + 2k\pi$$

Cộng $\dfrac{\pi}{3}$ vào cả ba vế: $-\dfrac{\pi}{6} + 2k\pi \le 2x \le \dfrac{5\pi}{6} + 2k\pi$. Chia cho $2$:

$$-\frac{\pi}{12} + k\pi \le x \le \frac{5\pi}{12} + k\pi$$

Vậy khoảng đồng biến là $\left[-\dfrac{\pi}{12} + k\pi,\ \dfrac{5\pi}{12} + k\pi\right]$, $k \in \mathbb{Z}$.
\end{dapan}

\begin{dapan}[12 — Đáp án D]
Phương trình $\cos 2x = 0$ khi $2x = \dfrac{\pi}{2} + k\pi$, suy ra $x = \dfrac{\pi}{4} + \dfrac{k\pi}{2}$ với $k \in \mathbb{Z}$.

\textbf{Chú ý.} Chia cả hai vế cho $2$ thì $k\pi$ cũng bị chia, thành $\dfrac{k\pi}{2}$. Quên chia là sai.
\end{dapan}

\begin{dapan}[13 — Đáp án D]
Vì $-1 \le \sin\left(x - \dfrac{\pi}{6}\right) \le 1$ nên $-5 \le 5\sin\left(x - \dfrac{\pi}{6}\right) \le 5$. Giá trị nhỏ nhất là $-5$.
\end{dapan}

\subsection*{Nhóm B — Công thức lượng giác}

\begin{dapan}[14 — Đáp án C]
Dùng hệ thức $\sin^2 x = 1 - \cos^2 x = (1 - \cos x)(1 + \cos x)$. Thay vào:

$$\frac{\sin^2 x}{1 - \cos x} = \frac{(1 - \cos x)(1 + \cos x)}{1 - \cos x} = 1 + \cos x$$

\textbf{Cách nhanh.} Nhận ra tử số là hiệu hai bình phương và rút gọn ngay, không cần biến đổi dài.
\end{dapan}

\begin{dapan}[15 — Đáp án D]
Công thức cộng của $\sin$: $\sin(\alpha + \beta) = \sin\alpha\cos\beta + \cos\alpha\sin\beta$.
\end{dapan}

\begin{dapan}[16 — Đáp án D]
Dùng cung hơn kém $180^\circ$: $\cos 210^\circ = \cos(180^\circ + 30^\circ) = -\cos 30^\circ = -\dfrac{\sqrt{3}}{2}$.
\end{dapan}

\begin{dapan}[17 — Đáp án B]
Dùng công thức cung bù: $\sin(180^\circ - \alpha) = \sin\alpha = \dfrac{1}{2}$.
\end{dapan}

\begin{dapan}[18 — Đáp án D]
Chia cả tử và mẫu cho $\cos\alpha$:

$$\frac{\sin\alpha - \cos\alpha}{\sin\alpha + \cos\alpha} = \frac{\tan\alpha - 1}{\tan\alpha + 1} = \frac{2 - 1}{2 + 1} = \frac{1}{3}$$
\end{dapan}

\begin{dapan}[19 — Đáp án B]
Từ $\tan\alpha = \sqrt{3}$ và $\tan\dfrac{\pi}{3} = \sqrt{3}$, suy ra $\alpha = \dfrac{\pi}{3}$ (một nghiệm đại diện).
\end{dapan}

\subsection*{Nhóm C — Hàm số bậc hai, mũ, logarit}

\begin{dapan}[20 — Đáp án C]
Từ $f(16) = 2$ và $f(x) = \log_a x$ được $\log_a 16 = 2$, tức $a^2 = 16$.

Vì $a > 0$ và $a \neq 1$ nên $a = 4$.
\end{dapan}

\begin{dapan}[21 — Đáp án A]
Nguyên hàm: $f(x) = \int (2x - 1)\,dx = x^2 - x + C$.

Điều kiện $f(0) = 1$ cho $C = 1$. Vậy $f(x) = x^2 - x + 1$.
\end{dapan}

\begin{dapan}[22 — Đáp án B]
Đạo hàm từng hạng tử: $\left(x\right)' = 1$ và $\left(\dfrac{1}{2x^2}\right)' = \left(\dfrac{1}{2}x^{-2}\right)' = -x^{-3} = -\dfrac{1}{x^3}$.

Vậy $f'(x) = 1 - \dfrac{1}{x^3}$.
\end{dapan}

\begin{dapan}[23 — Đáp án B]
Từ $\log_7 x = 2$ suy ra $x = 7^2 = 49$.
\end{dapan}

\subsection*{Nhóm D — Ôn tập tích lũy}

\begin{dapan}[24 — Đáp án D]
Giải $x^2 - 5x + 6 \le 0$: hai nghiệm $x = 2$, $x = 3$, lấy trong khoảng nên $A = \{2, 3\}$.

Giải $3x - 1 > 5$ được $x > 2$. Vì $B \subseteq \mathbb{N}$ nên $B = \{3, 4, 5, \ldots\}$.

Hợp hai tập: $A \cup B = \{2, 3, 4, 5, \ldots\} = \{x \in \mathbb{N} \mid x \ge 2\}$.
\end{dapan}

\begin{dapan}[25 — Đáp án A]
Công sai: $a_5 = a_1 + 4d$ cho $11 = 3 + 4d$, suy ra $d = 2$.

Tổng $7$ số hạng: $S_7 = \dfrac{7(a_1 + a_7)}{2}$ với $a_7 = 3 + 6 \cdot 2 = 15$. Vậy $S_7 = \dfrac{7 \cdot 18}{2} = 63$.
\end{dapan}

\subsection*{Nhóm D — Ôn tập tích lũy}

\begin{dapan}[1 — Đáp án A]
$\alpha$ ở phần tư thứ ba nên $\cos\alpha < 0$. Từ $\sin^2\alpha + \cos^2\alpha = 1$:

$$\cos\alpha = -\sqrt{1 - \frac{1}{4}} = -\frac{\sqrt{3}}{2}$$
\end{dapan}

\begin{dapan}[2 — Đáp án D]
$\tan\alpha = 1$ nên $\alpha$ có $|\sin\alpha| = |\cos\alpha|$. Vì $\alpha$ ở phần tư thứ ba nên $\sin\alpha < 0$, suy ra $\sin\alpha = -\dfrac{\sqrt{2}}{2}$.
\end{dapan}

\begin{dapan}[3 — Đáp án B]
Công thức cộng: $\tan\left(\alpha + \dfrac{\pi}{4}\right) = \dfrac{\tan\alpha + 1}{1 - \tan\alpha} = \dfrac{2 + 1}{1 - 2} = -3$.
\end{dapan}

\begin{dapan}[4 — Đáp án C]
$\cos$ là hàm chẵn nên $\cos\left(-\dfrac{\pi}{6}\right) = \cos\dfrac{\pi}{6} = \dfrac{\sqrt{3}}{2}$.
\end{dapan}

\begin{dapan}[5 — Đáp án D]
Đạo hàm hàm hợp: $f'(x) = \cos\left(2x + \dfrac{\pi}{4}\right) \cdot 2 = 2\cos\left(2x + \dfrac{\pi}{4}\right)$.
\end{dapan}

\begin{dapan}[6 — Đáp án C]
Gọi $R$ là bán kính đường tròn. Cung có số đo $2$ radian chắn dây cung độ dài $2$, nên nửa dây cung bằng $1$ và nửa góc ở tâm bằng $1$ radian.

Trong tam giác vuông: $\sin 1 = \dfrac{1}{R}$, suy ra $R = \dfrac{1}{\sin 1}$. Độ dài cung là $l = 2R = \dfrac{2}{\sin 1}$.
\end{dapan}

\begin{dapan}[7 — Đáp án A]
Với $a = 1$, $b = -1$: giá trị lớn nhất của $a\sin x + b\cos x$ là $\sqrt{a^2 + b^2} = \sqrt{2}$.
\end{dapan}

\begin{dapan}[8 — Đáp án A]
Dùng công thức nhân đôi: $\sin 2x\cos 2x = \dfrac{1}{2}\sin 4x$. Hệ số $\omega = 4$ nên chu kỳ là $\dfrac{2\pi}{4} = \dfrac{\pi}{2}$.
\end{dapan}

\begin{dapan}[9 — Đáp án A]
Biểu thức $-2\cos x$ lớn nhất bằng $2$ (khi $\cos x = -1$). Vậy $y = 3 - 2\cos x$ lớn nhất bằng $3 + 2 = 5$.

\textbf{Bẫy.} Dấu trừ trước biên độ làm đảo chiều: giá trị lớn nhất đạt khi $\cos x$ \textbf{nhỏ nhất}, không phải khi $\cos x = 1$.
\end{dapan}

\begin{dapan}[10 — Đáp án B]
Tính $f(-x) = (-x)\sin(-x) = (-x)(-\sin x) = x\sin x = f(x)$. Vậy hàm số là hàm chẵn.
\end{dapan}

\begin{dapan}[11 — Đáp án B]
Trên $[0, 2\pi]$, phương trình $\sin x = \dfrac{1}{2}$ có hai nghiệm $x = \dfrac{\pi}{6}$ và $x = \dfrac{5\pi}{6}$. Số nghiệm là $2$.
\end{dapan}

\begin{dapan}[12 — Đáp án B]
Đặt $t = \sin x$ với $t \in [-1, 1]$. Hàm $y = 2t^2 - 3t + 1$ có đỉnh $t_0 = \dfrac{3}{4}$, thuộc đoạn $[-1, 1]$.

Giá trị nhỏ nhất tại $t_0 = \dfrac{3}{4}$: $y = 2 \cdot \dfrac{9}{16} - \dfrac{9}{4} + 1 = -\dfrac{1}{8}$. Giá trị lớn nhất tại $t = -1$: $y = 2 + 3 + 1 = 6$.

Tập giá trị là $\left[-\dfrac{1}{8}, 6\right]$.
\end{dapan}

\begin{dapan}[13 — Đáp án C]
Đặt $t = \cos x$ với $t \in [-1, 1]$: $2t^2 - 3t + 1 = 0$ cho $t = 1$ hoặc $t = \dfrac{1}{2}$.

Với $t = 1$: $\cos x = 1$ cho $x = 0$ và $x = 2\pi$, tức $2$ nghiệm trên $[0, 2\pi]$.

Với $t = \dfrac{1}{2}$: $\cos x = \dfrac{1}{2}$ cho $x = \dfrac{\pi}{3}$ và $x = \dfrac{5\pi}{3}$, tức $2$ nghiệm. Tổng cộng $4$ nghiệm, nhưng $x = 0$ và $x = 2\pi$ trùng nhau khi đếm trên vòng tròn nên đáp án là $3$.
\end{dapan}

\begin{dapan}[14 — Đáp án C]
Đặt $t = \sin x$ với $t \in [-1, 1]$. Hàm $y = t^2 - t + 1$ có đỉnh $t_0 = \dfrac{1}{2}$, thuộc đoạn $[-1, 1]$.

Giá trị nhỏ nhất tại $t_0 = \dfrac{1}{2}$: $y = \dfrac{1}{4} - \dfrac{1}{2} + 1 = \dfrac{3}{4}$.
\end{dapan}

\begin{dapan}[15 — Đáp án D]
Trên khoảng $(0, \pi)$, phương trình $\tan x = -1$ có nghiệm $x = \dfrac{3\pi}{4}$.

\textbf{Chú ý.} Họ nghiệm tổng quát là $x = -\dfrac{\pi}{4} + k\pi$. Với $k = 1$ được $x = \dfrac{3\pi}{4}$, nằm trong $(0, \pi)$.
\end{dapan}

\clearpage
\sectiondautrang{Tổng kết lỗi thường gặp}
\addcontentsline{toc}{section}{Tổng kết lỗi thường gặp}

Bảng chẩn đoán dùng khi chữa bài và khi soạn đề kiểm tra. Cột ``Câu liên quan'' trỏ tới đề bài tập về nhà.

\bang{
\begin{tabular}{@{}p{5.6cm}p{4.2cm}p{5.2cm}@{}}
\toprule
\textbf{Lỗi} & \textbf{Câu liên quan} & \textbf{Cách cảnh báo học sinh} \\
\midrule
Nhầm chu kỳ $\tan$ là $2\pi$ & Câu 1, 20 phần bắt buộc & Riêng $\tan$ thì chia $\pi$, hai hàm kia chia $2\pi$ \\
Quên dấu trừ trước biên độ & Câu 9, 21 phần tự chọn & Viết $y = 3 - 2\cos x$, hỏi: lớn nhất là 1 hay 5? \\
Bỏ điều kiện $t \in [-1, 1]$ khi đặt ẩn phụ & Câu 10 phần bắt buộc & Giải ra $t = 2$ thì vô nghiệm, không phải $\cos x = 2$ \\
Nhầm số họ nghiệm của phương trình $\sin$ & Câu 7, 8 phần bắt buộc & $\sin$ có hai họ, $\cos$ và $\tan$ có một họ \\
Chia sai hằng số $k$ khi giải $\cos 2x = 0$ & Câu 12 phần bắt buộc & Chia hai vế cho 2 thì $k\pi$ thành $\dfrac{k\pi}{2}$ \\
Lẫn $\sin$ và $\cos$ khi dùng cung phụ & Câu 6 phần bắt buộc & $\cos\left(x - \dfrac{\pi}{2}\right) = \sin x$, dấu trừ cho dấu cộng \\
Không kiểm tra khoảng nghiệm của đề & Câu 9, 10 phần bắt buộc & Đọc kỹ đề có đoạn $[0, 2\pi]$ hay không \\
Quên đổi dấu khi dùng $\cos$ cung hơn kém $180^\circ$ & Câu 13 phần tự chọn & $\cos 210^\circ = -\cos 30^\circ$, có dấu trừ \\
\bottomrule
\end{tabular}}

% Không dùng \vfill: nó đẩy chữ xuống sát đáy trang,
% nơi banner letterhead chiếm chỗ (281-297mm) và sẽ che mất chữ.
\vspace{0.8cm}
\begin{center}
{\footnotesize\color{cae_blue} Tài liệu nội bộ, CAE SHANGHAI}
\end{center}

\end{document}
